\section{Feedback Quality Model}

\subsection{Six-Criterion Rubric}
Examiner Coach evaluates feedback using six criteria:
\begin{enumerate}
    \item specific observed behavior named;
    \item contextual feedback;
    \item objective and non-evaluative tone;
    \item explicitly mentioned strength;
    \item changeable area for improvement;
    \item concrete improvement plan or next step.
\end{enumerate}

These criteria were selected to represent feedback as an educationally useful
communication act. The first criterion asks whether the examiner names a
specific observed behavior, because learners can only improve reliably when they
know which action is being discussed. The second criterion checks whether the
comment is connected to the current OSCE station rather than being a generic
statement that could apply to any student. The third criterion addresses tone:
the system rewards descriptive, behavior-focused wording and penalizes
unsupported personal judgment.

The fourth and fifth criteria ensure that the feedback is balanced but not
superficial. A strength should be explicitly named, and an improvement area
should describe a changeable behavior rather than a fixed personal trait. The
sixth criterion captures actionability. It is not enough to tell a student that
something was missing; high-quality feedback should include a concrete next
step, such as a structure to follow, a phrase to use, or a specific clinical
sequence to practice.

\subsection{Scoring and Bilingual Output}
Each criterion is scored from 0 to 100 using predefined scoring anchors. The
system computes an overall score as the average of the six criterion scores and
counts criteria as met when they meet or exceed the default threshold
(\(\geq 70\)). Evaluation and coaching results can be resolved into German or
English without regenerating the whole evaluation.

The scoring model is intentionally simple. It does not claim to measure feedback
quality as an objective natural quantity; instead, it operationalizes feedback
quality for formative training. A score below 70 signals that a criterion is
absent or only weakly expressed, while a score of 70 or above indicates that the
criterion is present at a usable level. This makes the result easy to interpret
and allows the coaching component to focus on the weakest criteria.

Bilingual output is relevant because OSCE examiner training at TU Dresden may
involve German-language feedback while the technical and educational literature
is often available in English as in our case. The backend keeps evaluation content in
a structured form and resolves labels, summaries, and suggestions into the
requested language. This avoids regenerating the whole evaluation solely to
switch the display language.
